۱. backend/db.php — اتصال به دیتابیس
php
<?php
$host    = "localhost";
$db_user = "root";
$db_pass = "";
$db_name = "shop_db";

$conn = new mysqli($host, $db_user, $db_pass, $db_name);

if ($conn->connect_error) {
    die("خطا در اتصال به دیتابیس: " . $conn->connect_error);
}

$conn->set_charset("utf8mb4");
?>


این فایل چیکار می‌کنه؟
یه بار برای همیشه اطلاعات اتصال به MySQL رو تعریف می‌کنه و واقعاً وصل می‌شه. از این به بعد هر فایل PHP دیگه‌ای که بخواد با دیتابیس کار کنه (نمایش محصول، ثبت سفارش، پنل ادمین و...) کافیه بالاش بنویسه require '../backend/db.php'; و بعدش متغیر $conn رو داره — دیگه لازم نیست دوباره کد اتصال رو بنویسه. این یعنی اگه یه روز پسورد دیتابیس عوض شد، فقط همین یه فایل رو عوض می‌کنی، نه همه‌ی فایل‌ها رو.

خط set_charset("utf8mb4") هم فقط برای اینه که فارسی درست ذخیره و خونده بشه، بدونش ممکنه متن‌ها خراب نشون داده بشن.

۲. frontend/index.php — بالای فایل (قبل از <!DOCTYPE html>)
php
<?php
require '../backend/db.php';

$result = $conn->query("SELECT * FROM products");
?>

چیکار می‌کنه؟

خط اول → db.php رو صدا می‌زنه، یعنی الان $conn (اتصال به دیتابیس) در دسترسمونه
$conn->query("SELECT * FROM products") → یه دستور SQL می‌فرسته به MySQL: «همه ردیف‌های جدول products رو بده». نتیجه (همه ۷ محصول، هنوز خام و پردازش‌نشده) توی متغیر $result ذخیره می‌شه

اینجا فقط دیتا رو گرفتیم، هنوز نمایشش ندادیم — نمایش دادنش کار قسمت بعدیه.

۳. frontend/index.php — داخل بخش ProductsGrid
php
<div class="ProductsGrid">

   <?php while ($product = $result->fetch_assoc()): ?>
   <div class="ProductCard">
      <img class="ProductImg" src="<?php echo htmlspecialchars($product['image']); ?>" alt="<?php echo htmlspecialchars($product['name']); ?>">
      <p class="ProductName"><?php echo htmlspecialchars($product['name']); ?></p>
      <p class="ProductPrice"><?php echo number_format($product['price']); ?> تومان</p>
      <button class="AddToCartBtn">افزودن به سبد خرید</button>
   </div>
   <?php endwhile; ?>

</div>

چیکار می‌کنه؟ خط به خط:

<?php while ($product = $result->fetch_assoc()): ?> — یه حلقه شروع می‌شه. هر بار که اجرا می‌شه، fetch_assoc() یه ردیف از $result رو برمی‌داره و می‌ذاره توی $product (مثلاً اولین بار $product می‌شه محصول id=1، بار بعد id=2، و همینطور تا آخرین محصول). وقتی دیگه ردیفی نمونده، fetch_assoc() مقدار false برمی‌گردونه و حلقه خودش تموم می‌شه.
هرچی بین while و endwhile هست (یعنی خود <div class="ProductCard">) برای هر محصول یه بار کامل تکرار می‌شه. یعنی با ۷ محصول، این div واقعاً ۷ بار توی صفحه چاپ می‌شه.
$product['image'], $product['name'], $product['price'] — اینا دقیقاً همون اسم ستون‌های جدول products توی دیتابیس هستن. یعنی $product['name'] = محتوای همون سلولی که توی phpMyAdmin توی ستون name دیدی.
htmlspecialchars(...) — دور هر متنی که از دیتابیس میاد و توی HTML چاپش می‌کنیم می‌ذاریمش. کارش اینه که اگه یه‌وقت توی اسم محصول کاراکترهای خاص HTML بود (مثل < یا >)، اونا رو بی‌خطر می‌کنه که صفحه خراب نشه یا کسی نتونه کد مخرب تزریق کنه. یه عادت همیشگیه، نه چیز پیچیده‌ای.
number_format($product['price']) — عدد خام 450000 رو تبدیل می‌کنه به رشته 450,000 (کاما اضافه می‌کنه برای خوانایی بهتر).


--------------------------------------------------------------------------------